% #############################################################################
% This is Appendix B
% !TEX root = ../main.tex
% #############################################################################
\chapter{IR Specification}
\label{chapter:appendixIR}

This appendix specifies the syntax and the full instruction set of the \acf{IR} developed during this work. \Cref{tab:IRinstructionset} lists all instructions of the \ac{IR}, along with a brief description of each instruction. The full BNF syntax of \ac{IR} programs is provided in the \Cref{lst:IRsyntax}, containing the concrete syntax of all arguments used in the instructions.

In the instruction set specification, for convenience, the remote continuation referenced by a continuation \texttt{c} is denoted as \texttt{c.r}.

\setlength{\tabcolsep}{0pt}
\begin{table}
\centering
\caption{IR instruction set}
\label{tab:IRinstructionset}
\begin{tabular}{p{0.385\linewidth} p{0.61\linewidth}}
\hline
\textbf{Instruction} & \textbf{Description} \\
\hline
% \endfirsthead

\texttt{initContinuation(c, lab, loc)} &
Initializes continuation \texttt{c}, with return address \texttt{lab}, writing location \texttt{loc} and remote continuation \texttt{c}. \\

\texttt{bindContinuation(src, c, loc)} &
Binds the continuation stored at \texttt{src} to continuation \texttt{c}, setting the writing location of \texttt{c.r} to \texttt{loc}. \\

\texttt{continue(c, lab)} &
Transfers control to the return address of continuation \texttt{c}. Sets the return address of the remote continuation of \texttt{c} to \texttt{lab}. \\

\texttt{finish(c)} &
Transfers control to the return address of continuation \texttt{c}. Sets the remote continuation of the remote continuation of \texttt{c} to itself. \\

\texttt{forward(c1, c2)} &
Copies \texttt{c1} to \texttt{c2.r} and \texttt{c2} to \texttt{c1.r}, linking  \texttt{c1.r} and \texttt{c2.r}. Transfers control to the return address of \texttt{c2}. \\

\texttt{pushTask(lab)} &
Pushes a new task with code address \texttt{lab} onto the task queue. \\

\texttt{popTask()} &
Pops a task from the task queue and transfers control to its code address. \\

\texttt{jump(lab)} &
Unconditional jump to label \texttt{lab}. \\

\texttt{branchTag(loc, case, ...)} &
Branches on the tag stored at \texttt{loc}. Takes as arguments a list of cases for each tag \texttt{i}. \\

\texttt{branchExpr(expr, if, else)} &
Branches on the result of a boolean expression. Takes as arguments a \texttt{if} and an \texttt{else} case. \\

\texttt{branchFlags(reqs, if, else)} &
Similar to \texttt{branchExpr}, but branches on the whether the flag requirements \texttt{reqs} are satisfied. \\

\texttt{callProcess(pid, pargs)} &
Calls process \texttt{pid} with type parameters, continuation arguments and data arguments specified in \texttt{pargs}. \\

\texttt{callExponential(loc, c, d)} &
Calls the exponential closure stored at \texttt{loc}, initializing a new continuation \texttt{c} as its argument, with writing location \texttt{d}. \\

\texttt{writeExponential(loc, pid, eargs)} &
Writes to \texttt{loc} a new exponential closure from process \texttt{pid}, with type parameters and data arguments specified in \texttt{eargs}. \\

\texttt{writeExpression(loc, expr)} &
Evaluates \texttt{expr} and writes the result to \texttt{loc}. \\

\texttt{writeScan(loc, sl)} &
Reads a slot \texttt{sl} from the standard input and writes it to \texttt{loc}. \\

\texttt{writeTag(loc, i)} &
Writes a tag \texttt{i} to location \texttt{loc}. \\

\texttt{writeContinuation(loc, c)} &
Writes a continuation reference \texttt{c} to location \texttt{loc}. \\

\texttt{writeCell(loc, st)} &
Writes a new cell with capacity for slot tree \texttt{st} to location \texttt{loc}. \\

\texttt{writeType(loc, st, tflg)} &
Writes a type parameter to \texttt{loc} described by a slot tree \texttt{st} and with flag requisites \texttt{tflg}.\\

\texttt{moveSlots(dst, src, st)} &
Moves the slot tree \texttt{st} from location \texttt{src} into \texttt{dst}.\\

\texttt{cloneSlots(dst, src, st)} &
Clones the slot tree \texttt{st} from location \texttt{src} into \texttt{dst}.\\

\texttt{dropSlots(loc, st)} &
Drops the slot tree \texttt{st} at location \texttt{loc}.\\

\texttt{lockCell(loc)} &
Locks the cell at \texttt{loc}. \\

\texttt{unlockCell(loc)} &
Unlocks the cell at \texttt{loc}. \\

\texttt{acquireCell(loc)} &
Increases the reference count of the cell at \texttt{loc}. \\

\texttt{releaseCell(loc)} &
Decreases the reference count of the cell at \texttt{loc}. If the reference count reaches zero, the cell is deallocated. \\

\texttt{launchThread(lab)} &
Launches a new thread, starting execution at label \texttt{lab}. \\

\texttt{sleep(int)} &
Sleeps for a number of milliseconds specified by the integer argument. \\

\texttt{print(expr)} &
Evaluates \texttt{expr} and prints the result to the standard output. \\

\texttt{printLine(expr)} &
Same as \texttt{print}, but adds a newline character after printing. \\

\texttt{panic(str)} &
Prints the string \texttt{str} to the standard error and aborts execution. \\

\hline
\end{tabular}
\end{table}

\begin{figure}
\begin{lstlisting}[language=BNF,mathescape=true,frame=single,basicstyle=\ttfamily\footnotesize,caption={BNF Grammar of the Intermediate Representation},captionpos=t,label={lst:IRsyntax}]
nl   ::= "\n"
pid  ::= [a-zA-Z_][a-zA-Z0-9_]*
lab  ::= [a-zA-Z_][a-zA-Z0-9_]*
int  ::= [0-9]+
str  ::= /* double-quoted string */
bool ::= "true" | "false"

/* type, continuation, data and drop bit identifiers */
t, t1, t2 ::= "t" int
c, c1, c2 ::= "c" int
d, d1, d2 ::= "d" int
e         ::= "e" int

/* slots, slot trees, slot tree combinations, offsets */
sl   ::= "int" | "bool" | "string" | "tag"  | "cont"
       | "exponential"  | "cell"   | "type" | t
st   ::= "[]" | "[" sl (";" sl)* (";" st)? "]"
       | "tag[" st ("|" st)* "]"
       | "flags" "<" reqs ">" "[" st "|" st "]"
sc   ::= "<" st ("+" st)* ">"
ofs  ::= st "$\sim$" st

/* locations */
locb          ::= c | d | "c" "(" loc ")"
loc, src, dst ::= locb ("." ofs)?

/* type flag requirements */
reqs ::= "no" | "yes" | "if" "(" flag ("," flag)* ")"
flag ::= t "." ("value" | "droppable" | "cloneable")

/* type, continuation and data arguments */
tflg ::= ("(" flag "=" reqs ")")*
targ ::= t "<-" loc | t "<-" st tflg
carg ::= c1 "<-" (loc | c2) ("+" ofs)?
darg ::= d "<-[" st "]" loc | d "=[" st "]" loc 

/* process and exponential arguments */
parg  ::= targ | carg | darg
pargs ::= (parg ("," parg)*)?
earg  ::= targ | darg
eargs ::= (earg ("," earg)*)?

/* branch cases */
case, if, else ::= lab ":" int

/* program structure */
prog  ::= proc+
proc  ::= pid ":" nl entry* block+
entry ::= "exit points:" int nl
        | "types:" (t+ | "none") nl
        | "continuations:" (c+ | "none") nl
        | "data" d ":" sc ("(" c ")")? nl
        | "drop" e ":" st "at" loc "(always)"? nl
block ::= lab ":" nl (instr nl)+
expr  ::= expr "+" expr | expr "-" expr | expr "*" expr
        | expr "/" expr | expr "%" expr 
        | expr "=" expr | expr ">" expr | expr "<" expr
        | expr "and" expr | expr "or" expr | "not" expr
        | "(" expr ")" | int | str | bool
        | ("move" | "clone") "(" loc "," sl ")"
instr ::= /* consult instruction set table */
\end{lstlisting}
\end{figure}
